\section{Experimental Evaluation}
\label{sec:experiment}

\subsection{Experimental Setup}

We evaluate the Rust prototype corresponding to the construction in
Section~\ref{sec:construction}.  The implementation exposes the same top-level
interfaces as the construction: \textsf{Setup}, \textsf{SignEncap},
\textsf{Combine}, \textsf{Open}, \textsf{Verify}, and \textsf{Trace}.  The
baseline is the unmodified Gargos-style threshold Schnorr path.  The full
construction keeps the Gargos three-round signing flow, challenge generation,
response computation, aggregation rule, and final Schnorr verification equation
unchanged, while adding the auxiliary handle layer around finalized response
shares.

The experiments were run on macOS 14.5 on an Intel Core i5-1038NG7 MacBook Pro
with 16 GB RAM, 4 physical cores, and 8 logical cores.  The code was compiled in
Rust release mode using rustc 1.92.0 and cargo 1.92.0.  The group backend is
Ristretto over curve25519-dalek, with SHA-512 used for the domain-separated hash
functions in the prototype.  Unless otherwise stated, end-to-end measurements
use 100 independent trials after 10 warmup trials and report average, median,
standard deviation, and p95 latency.  Operations below timer precision are
measured by repeated execution.  Setup and key generation are not included in
online signing time.

The HTLP, IBE, and binding-proof backends used here are local
interface-compatible prototype components.  They implement the APIs required by
the construction and allow integration testing, but they are not production HTLP,
IBE, or general-purpose NIZK implementations.  Consequently, the absolute delay,
proof size, and encryption timings should be read as prototype measurements; the
protocol-level structure and object counts are the parts evaluated here.

\subsection{End-to-End Performance}

Table~\ref{tab:e2e-performance} compares the unmodified Gargos baseline with the
full construction.  The parameter choice for this table is $n=|S|$ and
$t=|S|-1$, so all listed signers participate.  This isolates scaling in the
active signer set.  The full construction preserves the same external Schnorr
verification equation, but adds signer-local auxiliary generation and public
record validation.

\begin{table}[t]
\centering
\caption{End-to-end runtime for Gargos baseline and the full construction.}
\label{tab:e2e-performance}
\small
\begin{tabular}{llrrrr}
\toprule
Mode & $|S|$ & Sign & Combine & Verify & Total \\
\midrule
Gargos & 4 & 2.563 & 1.845 & 0.067 & 4.542 \\
Full & 4 & 4.732 & 0.450 & 0.197 & 5.905 \\
Gargos & 8 & 5.052 & 3.631 & 0.068 & 9.380 \\
Full & 8 & 9.926 & 0.761 & 0.257 & 11.603 \\
Gargos & 16 & 10.569 & 8.074 & 0.068 & 18.365 \\
Full & 16 & 19.346 & 1.410 & 0.375 & 22.082 \\
Gargos & 32 & 21.359 & 14.993 & 0.067 & 36.353 \\
Full & 32 & 38.634 & 2.782 & 0.627 & 43.728 \\
Gargos & 64 & 43.557 & 30.165 & 0.068 & 73.671 \\
Full & 64 & 78.391 & 5.406 & 1.135 & 87.760 \\
\bottomrule
\end{tabular}
\end{table}

The full construction adds measurable computation but does not add an extra
interactive signing round.  At $|S|=64$, median end-to-end time increases from
73.671 ms to 87.760 ms, an increase of 14.088 ms (1.19x). This is a modest computational overhead, but it is not negligible.

\subsection{Component Breakdown}

Table~\ref{tab:phase-breakdown} separates the main benchmark phases at
$|S|=32$.  The rows distinguish baseline signing, signer-local auxiliary handle
generation, normal combining, delayed opening, public verification, tracing, and
end-to-end execution.  The delayed opening path is split into puzzle aggregation
(\textsf{PEval}) and puzzle solving (\textsf{PSolve}).

\begin{table}[t]
\centering
\caption{Phase-level runtime breakdown at $|S|=32$.}
\label{tab:phase-breakdown}
\small
\begin{tabular}{lrrrr}
\toprule
Phase & Avg. & Median & Std. dev. & P95 \\
\midrule
Gargos baseline signing & 36.667 & 36.382 & 1.052 & 38.631 \\
Auxiliary handle generation per signer & 0.080 & 0.072 & 0.051 & 0.102 \\
Commitment generation & 0.063 & 0.058 & 0.047 & 0.078 \\
Two HTLP puzzle generations & 0.003 & 0.003 & 0.002 & 0.003 \\
IBE encryption & 0.002 & 0.001 & 0.014 & 0.002 \\
Binding-proof generation & 0.007 & 0.007 & 0.003 & 0.007 \\
Normal Combine & 41.974 & 41.828 & 1.244 & 44.593 \\
Delayed Open PEval & 39.832 & 39.508 & 1.425 & 42.463 \\
Delayed Open PSolve & 0.002 & 0.002 & 0.000 & 0.002 \\
Public Verify & 42.825 & 42.534 & 1.305 & 44.955 \\
Trace & 42.477 & 42.319 & 0.919 & 44.401 \\
End-to-end baseline & 36.792 & 36.609 & 0.948 & 38.446 \\
End-to-end full construction & 44.102 & 44.012 & 0.995 & 45.479 \\
\bottomrule
\end{tabular}
\end{table}

Table~\ref{tab:aux-components} drills down into the signer-local auxiliary
operations.  In this prototype, commitment generation is the largest single
auxiliary primitive cost.  The HTLP, IBE, nullifier, and prototype binding-proof
rows are lightweight because they are implemented by local prototype backends.
A production NIZK backend would change the binding-proof row and must be
measured separately.

\begin{table}[t]
\centering
\caption{Signer-local auxiliary component costs.}
\label{tab:aux-components}
\small
\begin{tabular}{lrrrr}
\toprule
Component & Avg. & Median & Std. dev. & P95 \\
\midrule
Commitment generation & 0.090518 & 0.058664 & 0.279926 & 0.125066 \\
HTLP generation & 0.001490 & 0.001464 & 0.000318 & 0.001599 \\
IBE encryption & 0.001151 & 0.001158 & 0.000078 & 0.001198 \\
NIZK proof generation & 0.008361 & 0.007076 & 0.018512 & 0.008523 \\
Nullifier computation & 0.000390 & 0.000389 & 0.000007 & 0.000393 \\
\bottomrule
\end{tabular}
\end{table}

\subsection{Delayed Recovery and Tracing}

Table~\ref{tab:htlp-delay} reports the prototype HTLP solving time for different
sequential-work parameters.  The parameter is the number of sequential hash
iterations used by the local benchmark backend.  Changing it affects delayed
release latency, but not the construction interface, the canonical record, or
the verification equations.

\begin{table}[t]
\centering
\caption{Prototype HTLP solving time under different delay parameters.}
\label{tab:htlp-delay}
\small
\begin{tabular}{rrrr}
\toprule
Delay iterations & Median & Std. dev. & P95 \\
\midrule
0 & 0.000 & 0.000 & 0.000 \\
1000 & 0.343 & 0.085 & 0.464 \\
5000 & 1.721 & 0.264 & 2.530 \\
10000 & 3.619 & 0.621 & 5.398 \\
20000 & 7.303 & 0.585 & 8.431 \\
\bottomrule
\end{tabular}
\end{table}

The \textsf{Open} and \textsf{Trace} rows in Table~\ref{tab:e2e-performance} show
that both operations scale with the number of public handles, because each
handle must be checked and, for tracing, decrypted under the message-dependent
identity.  The delayed path verifies the same aggregate commitment and Schnorr
equation as normal release; no check is removed from the recovery path.

\subsection{Communication Overhead}

Table~\ref{tab:communication-total} reports serialized communication sizes for
$|S|=32$.  The full construction includes the baseline Gargos transcript, all
public handles in the canonical record, all private final-round packages, and
the final auxiliary signature package.  Aggregate HTLP objects are also shown as
a subcomponent of the final package, so they are not added a second time when
computing the total.

\begin{table}[t]
\centering
\caption{Serialized communication and storage components at $|S|=32$.}
\label{tab:communication-total}
\small
\begin{tabular}{lr}
\toprule
Item & Bytes \\
\midrule
Baseline Gargos transcript & 11756 \\
One public handle & 292 \\
Complete canonical public record & 9468 \\
One private final-round package & 128 \\
Aggregate HTLP objects & 128 \\
Final auxiliary signature package & 308 \\
Total full construction & 25628 \\
Additional bytes relative to Gargos & 13872 \\
Per-signer marginal overhead & 433 \\
\bottomrule
\end{tabular}
\end{table}

Table~\ref{tab:handle-composition} breaks down one public handle.  The increased
communication cost is explicit: at $|S|=32$, the full construction uses 25628
bytes, compared with 11756 bytes for the baseline transcript.  The additional
13872 bytes are linear in the number of participating signers.

\begin{table}[t]
\centering
\caption{Composition of one public auxiliary handle.}
\label{tab:handle-composition}
\small
\begin{tabular}{lr}
\toprule
Handle field & Bytes \\
\midrule
commitment & 32 \\
two HTLPs & 128 \\
IBE ciphertext & 68 \\
nullifier & 32 \\
binding proof & 32 \\
\bottomrule
\end{tabular}
\end{table}

\subsection{Scalability}

Table~\ref{tab:scalability-nt} varies the threshold configuration.  For each
pair $(n,t)$, the active set has minimum authorized size $\ell=t+1$.  The result
shows that runtime grows with the number of participating signers, because one
public handle and one private package are generated per active signer.

\begin{table}[t]
\centering
\caption{Scalability under different threshold parameters.}
\label{tab:scalability-nt}
\small
\begin{tabular}{rrrrrr}
\toprule
$n$ & $t$ & $\ell$ & Sign & Verify & Total \\
\midrule
10 & 2 & 3 & 3.380 & 0.181 & 4.330 \\
20 & 4 & 5 & 5.714 & 0.210 & 6.994 \\
32 & 8 & 9 & 11.773 & 0.273 & 13.744 \\
48 & 16 & 17 & 20.687 & 0.390 & 23.612 \\
\bottomrule
\end{tabular}
\end{table}

The data support near-linear growth in the tested range, with active-set size as
the main driver.  We do not claim asymptotic improvement over Gargos; the
construction adds per-signer transcript objects while preserving Gargos's
three-round interaction pattern.

\subsection{Correctness and Robustness Tests}

Table~\ref{tab:correctness-rates} reports 1000 independent sessions for each
active signer set.  The tests check public verification, delayed recovery,
tracing accuracy, and equality between normal and delayed release outputs.

\begin{table}[t]
\centering
\caption{Correctness rates over 1000 independent sessions.}
\label{tab:correctness-rates}
\small
\begin{tabular}{rrrrr}
\toprule
$|S|$ & Verify & Delayed recovery & Trace accuracy & Equal outputs \\
\midrule
4 & 100.0\% & 100.0\% & 100.0\% & 100.0\% \\
8 & 100.0\% & 100.0\% & 100.0\% & 100.0\% \\
16 & 100.0\% & 100.0\% & 100.0\% & 100.0\% \\
32 & 100.0\% & 100.0\% & 100.0\% & 100.0\% \\
64 & 100.0\% & 100.0\% & 100.0\% & 100.0\% \\
\bottomrule
\end{tabular}
\end{table}

Table~\ref{tab:robustness-rejection} reports rejection tests over 1000 sessions
at $|S|=8$.  Each row modifies one object or consistency condition and then
checks that verification or tracing rejects the resulting transcript.

\begin{table}[t]
\centering
\caption{Robustness rejection tests over 1000 sessions.}
\label{tab:robustness-rejection}
\small
\begin{tabular}{lrr}
\toprule
Modified object or condition & Rejected & Rate \\
\midrule
modified commitment & 1000/1000 & 100.0\% \\
modified puzzle & 1000/1000 & 100.0\% \\
modified ciphertext & 1000/1000 & 100.0\% \\
modified nullifier & 1000/1000 & 100.0\% \\
modified proof & 1000/1000 & 100.0\% \\
modified message & 1000/1000 & 100.0\% \\
modified session digest & 1000/1000 & 100.0\% \\
modified aggregate response & 1000/1000 & 100.0\% \\
duplicate nullifier & 1000/1000 & 100.0\% \\
missing handle & 1000/1000 & 100.0\% \\
cross-session handle & 1000/1000 & 100.0\% \\
wrong tracing key & 1000/1000 & 100.0\% \\
\bottomrule
\end{tabular}
\end{table}

\subsection{Comparison and Limitations}

Table~\ref{tab:functional-comparison} compares functionality with Gargos and
TiMTAPS.  We do not report absolute TiMTAPS runtime numbers because the current
prototype does not reproduce TiMTAPS on the same primitives, curve, and machine.
A direct timing comparison would therefore be misleading.  The comparison is
restricted to functionality and protocol structure.

\begin{table}[t]
\centering
\caption{Functional comparison with related schemes.}
\label{tab:functional-comparison}
\small
\begin{tabular}{lccccc}
\toprule
Scheme & Adaptive security & 3-round signing & Timed combining & Verifiable combining & Message-dependent tracing \\
\midrule
Gargos & Yes & Yes & No & No & No \\
TiMTAPS & Not Gargos-based & No & Yes & Yes & Yes \\
Ours & Inherits Gargos layer & Yes & Yes & Yes & Yes \\
\bottomrule
\end{tabular}
\end{table}

The results show the tradeoff introduced by the auxiliary layer.  The protocol
keeps Gargos's three-round signing semantics and standard Schnorr verification,
but adds signer-local computation and a larger transcript.  The computational
overhead is modest in this prototype, while communication increases to about
2.18x the baseline at $|S|=32$.  These costs are the price of delayed recovery,
public aggregation consistency, and message-dependent tracing.  Since HTLP, IBE,
and NIZK are prototype backends here, production deployments must remeasure the
corresponding rows with the chosen concrete primitives.
